\section{The colour shortcut, measured}
\label{sec:shortcut}

Our report found that colour jitter added 15--21 points to every structural pipeline (\cref{app:report}) and
attributed this to a colour-histogram shortcut \citep{chen2020simclr,geirhos2020shortcut}: without colour
distortion, two crops of an image can be matched by their colour statistics alone. The reviewer asked for
evidence of the mechanism rather than of its effect. We trained \simclr on all images with four augmentation
pipelines that form the three tiers of the report and measured how much each encoder relies on colour
(\cref{tab:color}):
\begin{itemize}[leftmargin=*]
\item \textbf{Gray test images.} We convert the test images to grayscale and apply the unchanged linear probe.
An encoder that matches views by colour should suffer most.
\item \textbf{Histogram alignment.} Linear CKA \citep{kornblith2019cka} between the encoder features of the test
images and their joint RGB histograms ($8^3$ bins, no spatial information), next to the CKA with the class
labels; and the $R^2$ of a ridge regression that predicts the histogram from the features.
\end{itemize}

\begin{table}[t]
\centering
\small
\caption{Colour reliance of \simclr encoders trained with different augmentations (all images, one seed).
$\Delta$: change in linear accuracy when the test images are converted to grayscale. CKA(hist.): linear CKA
between features and colour histograms; CKA(label): between features and one-hot labels.}
\label{tab:color}
\begin{adjustbox}{max width=\linewidth}
\begin{tabular}{l l ccc ccc}
\toprule
 & & \multicolumn{3}{c}{accuracy (\%)} & \multicolumn{3}{c}{colour reliance} \\
\cmidrule(lr){3-5}\cmidrule(lr){6-8}
augmentation & tier & linear & gray test & $\Delta$ & CKA(hist.) & CKA(label) & hist.\ $R^2$ \\
\midrule
crop + flip & Tier 1 & 60.9 & 24.5 & -36.4 & 0.57 & 0.16 & 0.64 \\
crop + flip + jitter & Tier 2 & 75.9 & 38.2 & -37.6 & 0.47 & 0.28 & 0.48 \\
crop + flip + grayscale & Tier 2 & 77.1 & 75.5 & -1.6 & 0.30 & 0.34 & 0.33 \\
crop + flip + jitter + grayscale & Tier 3 & 83.8 & 82.7 & -1.1 & 0.22 & 0.47 & 0.25 \\
\midrule
colour histogram only ($8^3$ bins, no encoder) & -- & 37.7 & -- & -- & 1.00 & -- & 1.00 \\
\bottomrule
\end{tabular}

\end{adjustbox}
\end{table}

Without colour distortion (crop + flip, \augTwoLin\% linear accuracy) the encoder is a colour detector: its
features are more aligned with the colour histogram than with the class labels (CKA \augTwoCka against
\augTwoCkaLabel), the histogram can be read off them almost linearly, and converting the test images to gray
costs \augTwoGrayDrop points (to \augTwoGrayAcc\%), which is less than a linear classifier on the raw colour
histogram achieves on its own (\histProbe\%). This is the shortcut, measured directly. Each colour transform removes part of it.
Random grayscale lowers the histogram alignment most and, as expected, makes the encoder insensitive to gray
test images (a drop of \augSixGrayDrop points), while colour jitter alone leaves a large gray-image penalty (\augFiveGrayDrop points) because
it perturbs but rarely removes colour. Accuracy needs both (\augFiveLin and \augSixLin\% with one of them,
\augEightLin\% with both). With grayscale alone the encoder is about as aligned with the colour histogram as with
the labels (CKA \augSixCka vs.\ \augSixCkaLabel); with both transforms the label alignment is twice the histogram
alignment (\augEightCkaLabel vs.\ \augEightCka). The tiers of our report thus order the encoders by how much colour information they keep.

Labels make the shortcut much less important. Removing colour jitter from the full pipeline costs \simclr
\augJitterCost points (\augSixLin\% instead of \augEightLin\%), but label-only \supcon only \supconJitterCost
(\supconNoJitterLin\% instead of \supconLin\%): when two images of the same class are positives, matching colour
statistics is no longer enough to solve the task, so the labels themselves break the shortcut. This is a second
way in which labels buy speed: they substitute for part of the augmentation that a self-supervised method needs.

\paragraph{Geometry of the augmentation tiers.}
\citet{luthra2025approximately} predict that stronger invariances push the representation towards a collapsed
simplex ETF. \Cref{tab:etf} confirms this quantitatively on the test set: from crop + flip to the full pipeline,
the CDNV falls from \augTwoCdnv to \augEightCdnv and the mean deviation of the class-mean angles from those of a
simplex ETF halves (\augTwoEtf to \augEightEtf), with the single-transform tiers in between. The class means
are equiangular on average for every encoder (mean cosine close to $-1/9$); what changes is how evenly the
angles are spread around that value. The t-SNE plots of our report showed this qualitatively; the ETF
statistics turn it into a measurement.

\begin{table}[t]
\centering
\small
\caption{Neural-collapse statistics on the CIFAR-10 test set for the augmentation tiers (one seed). A simplex
ETF has pairwise cosines of $-1/9=-0.111$ between centred class means, ETF deviation 0 and equinorm CV 0.}
\label{tab:etf}
\begin{adjustbox}{max width=\linewidth}
\begin{tabular}{l l ccccc}
\toprule
augmentation & tier & linear & CDNV$\downarrow$ & ETF dev.$\downarrow$ & $\cos$ mean (target $-0.111$) & equinorm CV$\downarrow$ \\
\midrule
crop + flip & Tier 1 & 60.9 & 7.24 & 0.462 & -0.108 & 0.180 \\
crop + flip + jitter & Tier 2 & 75.9 & 2.78 & 0.350 & -0.099 & 0.230 \\
crop + flip + grayscale & Tier 2 & 77.1 & 2.22 & 0.314 & -0.098 & 0.243 \\
crop + flip + jitter + grayscale & Tier 3 & 83.8 & 1.23 & 0.226 & -0.100 & 0.224 \\
\bottomrule
\end{tabular}

\end{adjustbox}
\end{table}

\section{A resolution boundary, quantified}
\label{sec:ood}

The retrieval demo of our original project returned horses and deer when queried with a high-resolution
photograph of a German Shepherd. To check that this is a systematic effect and not an anecdote, we took all
\nOod validation images of Imagewoof \citep{howard2019imagewoof}, ten dog breeds from ImageNet
\citep{deng2009imagenet}, centre-cropped them and downsampled them to $32{\times}32$ with an antialiased
bicubic filter, as the demo does. We then measured how often each encoder's linear probe calls them ``dog''
and how many of their five nearest CIFAR-10 test images are dogs (\cref{tab:ood}).

\begin{table}[t]
\centering
\small
\caption{Resolution shift: ImageNet dog photographs ($\geq$320\,px) downsampled to 32\,px, compared with CIFAR-10
test dogs. Dog recall of the linear probe (95\% bootstrap CI) and retrieval precision@5 in the CIFAR-10 test
set (seed 0 encoders). CLIP (zero-shot, ViT-B/32) is an external reference trained on 400M image--text
pairs.}
\label{tab:ood}
\begin{adjustbox}{max width=\linewidth}
\begin{tabular}{l cc cc}
\toprule
 & \multicolumn{2}{c}{dog recall (\%)} & \multicolumn{2}{c}{retrieval P@5 (\%)} \\
\cmidrule(lr){2-3}\cmidrule(lr){4-5}
encoder & CIFAR-10 dogs & ImageNet dogs$\to$32px & CIFAR-10 dogs & ImageNet dogs$\to$32px \\
\midrule
SimCLR & 73.4 & 42.2 {\scriptsize[40.6, 43.7]} & 61.1 & 36.6 \\
SupCon (label-only) & 74.5 & 49.6 {\scriptsize[48.1, 51.2]} & 63.9 & 41.9 \\
Union & 77.6 & 47.0 {\scriptsize[45.4, 48.6]} & 66.1 & 42.0 \\
SimCLR, 19,400 steps & 77.6 & 46.8 {\scriptsize[45.2, 48.4]} & 68.4 & 44.6 \\
Supervised CE (all labels) & 90.8 & 64.7 {\scriptsize[63.3, 66.2]} & 87.8 & 61.8 \\
\midrule
CLIP ViT-B/32 zero-shot (reference) & 87.1 & 83.7 & -- & -- \\
\bottomrule
\end{tabular}

\end{adjustbox}
\end{table}

Every CIFAR-trained encoder recognises far fewer of the downsampled ImageNet dogs than of the CIFAR-10 dogs
(\simclr: \oodSimclrCifar\% $\to$ \oodSimclrOod\%), and the most frequent wrong answer is
\oodSimclrTopWrong. The drop is not specific to self-supervision: the supervised encoder trained with all labels
falls from \oodCeCifar\% to \oodCeOod\%. CLIP shows that most of the information survives the
downsampling (dog recall \clipResThirtyTwo\% at 32\,px vs.\ \clipResTwoTwoFour\% at 224\,px). The failure is
therefore a distribution shift between ImageNet photographs and CIFAR-10 images at the same resolution
(framing, pose, and the downsampling filter), which all of our encoders inherit from CIFAR-10. We had
attributed the demo failure to the lack of language supervision; the supervised baseline shows that the
training distribution, not the kind of supervision, is the main cause.
